\documentclass[11pt,a4paper]{article}

% ── Packages ─────────────────────────────────────────────────────────────────
\usepackage[T1]{fontenc}
\usepackage[utf8]{inputenc}
\usepackage[margin=2.5cm]{geometry}
\usepackage{amsmath,amssymb,amsthm}
\usepackage{booktabs}
\usepackage{xcolor}
\usepackage{listings}
\usepackage{hyperref}
\usepackage{microtype}
\usepackage{parskip}
\usepackage{titlesec}
\usepackage{enumitem}
\usepackage{caption}
\usepackage{float}
\usepackage{lmodern}

% ── Colour palette ────────────────────────────────────────────────────────────
\definecolor{codebg}{rgb}{0.97,0.97,0.97}
\definecolor{keyword}{rgb}{0.00,0.30,0.70}
\definecolor{comment}{rgb}{0.40,0.40,0.40}
\definecolor{string}{rgb}{0.60,0.10,0.10}
\definecolor{lineno}{rgb}{0.55,0.55,0.55}
\definecolor{headerbg}{rgb}{0.92,0.95,1.00}

% ── Code listing style ────────────────────────────────────────────────────────
\lstdefinestyle{cpp}{
  language=C++,
  backgroundcolor=\color{codebg},
  basicstyle=\ttfamily\small,
  keywordstyle=\color{keyword}\bfseries,
  commentstyle=\color{comment}\itshape,
  stringstyle=\color{string},
  numbers=left,
  numberstyle=\tiny\color{lineno},
  numbersep=6pt,
  frame=single,
  framerule=0.4pt,
  rulecolor=\color{black!20},
  breaklines=true,
  breakatwhitespace=false,
  tabsize=4,
  showspaces=false,
  showstringspaces=false,
  captionpos=b,
  morekeywords={uint64_t,int32_t,int64_t,uint8_t,ptrdiff_t,size_t,
                __builtin_popcountll,__device__,__global__,
                cuDoubleComplex,static_assert,constexpr,noexcept,nodiscard,
                make_cuDoubleComplex,__popcll,atomicAdd,
                shared_ptr,unique_ptr,vector,pragma},
}

\lstdefinestyle{pseudo}{
  backgroundcolor=\color{codebg},
  basicstyle=\ttfamily\small,
  frame=single,
  framerule=0.4pt,
  rulecolor=\color{black!20},
  breaklines=true,
  numbers=none,
  captionpos=b,
}

\lstset{style=cpp}

% ── Section formatting ────────────────────────────────────────────────────────
\titleformat{\section}{\large\bfseries}{{\thesection}}{1em}{}
\titleformat{\subsection}{\normalsize\bfseries}{{\thesubsection}}{1em}{}
\titleformat{\subsubsection}{\normalsize\itshape}{{\thesubsubsection}}{1em}{}

% ── Hyperref setup ────────────────────────────────────────────────────────────
\hypersetup{
  colorlinks=true,
  linkcolor=blue!60!black,
  citecolor=green!40!black,
  urlcolor=blue!70!black,
  pdftitle={Spatial Symmetry Workflow: Architecture and Optimizations},
  pdfauthor={QED Development},
}

% ── Math macros ──────────────────────────────────────────────────────────────
\newcommand{\G}{|\mathcal{G}|}          % |G|
\newcommand{\ket}[1]{|{#1}\rangle}
\newcommand{\rb}{r_b}
\newcommand{\chik}{\chi_k}

% ── Theorem environments ─────────────────────────────────────────────────────
\newtheorem{remark}{Remark}

% ─────────────────────────────────────────────────────────────────────────────
\begin{document}

\begin{center}
  {\LARGE\bfseries Spatial Symmetry Workflow}\\[4pt]
  {\large Architecture, Implementation, and Performance Optimizations}\\[8pt]
  {\normalsize QED Exact Diagonalization Codebase \quad$\cdot$\quad June 2026}
\end{center}

\vspace{8pt}
\hrule
\vspace{12pt}

\tableofcontents

\newpage

% =============================================================================
\section{Overview}
% =============================================================================

This document covers the complete call chain for the spatial symmetry workflow
in the QED exact diagonalization code, from the Python \texttt{qed.thermal()}
API down to the CUDA device kernel.  It documents every optimization applied in
June 2026, including the mathematical basis, implementation details, memory
analysis, and measured performance impact.

The primary entry point under study is
\begin{center}
\texttt{workflows\_thermal\_all\_sz\_streaming\_symmetry\_directory()}
\end{center}
which computes finite-temperature thermodynamics for a lattice Hamiltonian using
all $S_z$ sectors combined with spatial symmetry (translational + point-group),
running all resulting sectors in a single flat OpenMP parallel loop.

% =============================================================================
\section{Architecture: Two Code Paths}
\label{sec:arch}
% =============================================================================

The workflow splits at the matrix-vector product (matvec).  Setup is shared;
execution diverges into the \emph{rep path} (default) or the \emph{CSR path}
(fallback).

\subsection{Call chain summary}

\begin{lstlisting}[style=pseudo, caption={Top-level call chain from Python to C++/CUDA kernel}]
Python: qed.thermal(..., use_symmetry_if_available=True,
                        use_sz_if_conserved=True)
  |
  +-- workflows_thermal_all_sz_streaming_symmetry_directory()
            |                              [workflow_bindings.cpp]
            |
            +-- [SETUP] make_all_sz_sector_operators_tagged()
            |      |                       [make_operator.h]
            |      +-- build_all_sz_sector_operators()
            |                              [sector_set.h]
            |            +-- enumerate_full_orbit_reps()
            |            +-- Pass 1.5: compute_orbit_for_state()
            |            |                 [projector_chain.h]
            |            +-- SectorOperator::configureRepLazy()
            |                              [sector_basis.h]
            |
            +-- [HOT] #pragma omp parallel for  (156 sectors)
                        |
                        +-- CPU rep path  (default, ED_SYM_REP=1):
                        |    CpuMatVecBackend<RepSymmetryBasisPolicy>
                        |    +-- index_and_projection()
                        |            [rep_symmetry_basis_policy.h]
                        |
                        +-- GPU rep path  (allow_gpu=True, large dim):
                             DeviceRepSymmetryBasisPolicy
                             +-- index_and_projection()
                                     [device_basis_policy.cuh]
\end{lstlisting}

\subsection{Rep path vs.\ CSR path}

\begin{description}[leftmargin=2.5cm, labelwidth=2.3cm]
  \item[\textbf{Rep path}] Stores only orbit representatives and
    $1/\|\cdot\|$ norms.  The group action and projection phase are
    regenerated \emph{on the fly} inside \texttt{index\_and\_projection}.
    No orbit-CSR matrix is materialised.  Active when
    \texttt{ED\_SYM\_REP=1} (default) and \texttt{RepSectorData::usable()}
    returns true.

  \item[\textbf{CSR path}] Pre-bakes every orbit image and its
    symmetry-projected coefficient into a hash table (GPU) or sorted array
    (CPU) at setup time.  Higher setup cost; zero \texttt{apply\_perm} calls
    during the matvec.  Activated when \texttt{ED\_SYM\_REP=0} or when the
    sector cannot carry rep data (\emph{e.g.}, $S_z$ is not conserved).
\end{description}

\medskip
For $N = 12$ with $|\mathcal{G}| = 12$ (cyclic translation group of the
Heisenberg chain): $13$ $S_z$ sectors $\times$ $12$ irreps $= 156$ total
\texttt{SectorOperator} objects are built and solved in parallel.

% =============================================================================
\section{Setup Pipeline (CPU, runs once per call)}
\label{sec:setup}
% =============================================================================

\subsection{JSON loading and group decomposition}

\textbf{File:} \texttt{QED/include/ed/symmetry/symmetry\_group\_info.h}

\medskip
\texttt{SymmetryGroupInfo::loadFromDirectory()} reads three JSON files from the
\texttt{automorphism\_results/} directory:

\begin{enumerate}
  \item \texttt{max\_clique.json} $\to$ \texttt{max\_clique[g][site]}: the full
    group as a list of site-permutation arrays (one per group element $g$).
  \item \texttt{minimal\_generators.json} $\to$ generator list with orders.
  \item \texttt{sector\_metadata.json} $\to$
    \texttt{sectors[k].phase\_factors[]}: per-generator complex characters
    $\phi_{k,j}$ for each irrep $k$.
\end{enumerate}

After loading, \texttt{computePowerRepresentation()} performs a BFS on the
generator Cayley graph to express every $g \in \mathcal{G}$ as a product of
generator powers:
\begin{equation}
  g \;=\; \mathrm{gen}_0^{\,p_0} \circ \mathrm{gen}_1^{\,p_1} \circ \cdots
  \label{eq:power_rep}
\end{equation}
storing the exponents as \texttt{power\_representation[g][j] = $p_j$}.
This allows the sector character to be evaluated as
\begin{equation}
  \chi_k(g) \;=\; \prod_{j=0}^{n_{\mathrm{gen}}-1} \phi_{k,j}^{\,p_j}
  \label{eq:character}
\end{equation}
without storing one complex number per $(k,g)$ pair explicitly (only the
generator characters $\phi_{k,j}$ are stored).

\subsection{Orbit representative enumeration}

\textbf{File:} \texttt{QED/include/ed/symmetry/orbit\_reps.h}

\begin{lstlisting}[style=pseudo, caption={Pseudocode for full-space orbit enumeration}]
enumerate_full_orbit_reps(info, n_bits):
  all_reps = []
  for s = 0 .. 2^N - 1:
      r = min_{g in G}  apply_perm(s, g)
      if r == s:          # s is its own orbit representative
          all_reps.append(s)
  return sorted(all_reps)   # 2^N / |G| entries
\end{lstlisting}

The result is partitioned by Hamming weight (popcount) into
\texttt{reps\_by\_n\_up[0\ldots N]}: since lattice automorphisms preserve
particle number, every group image of a fixed-$S_z$ state has the same
$n_{\uparrow}$.

For $N = 12$, $|\mathcal{G}| = 12$: scans $4096$ states, yields
$\approx 342$ representatives split across $13$ $S_z$ sectors
($n_{\uparrow} = 0, 1, \ldots, 12$).

\subsection{Pass 1.5: \texttt{RepSectorData} build}
\label{sec:pass15}

\textbf{File:} \texttt{QED/include/ed/symmetry/sector\_set.h}, function
\texttt{build\_all\_sz\_sector\_operators()}

This is the critical setup loop.  For each pair $(n_{\uparrow}, k)$ where
$k$ is the irrep index, it:
\begin{enumerate}
  \item Calls \texttt{compute\_orbit\_for\_state()} (Section~\ref{sec:orbit})
    for every rep in \texttt{reps\_by\_n\_up[$n_\uparrow$]} to check whether
    that rep survives in irrep $k$ (i.e., the projected norm exceeds a
    threshold~$\epsilon$).
  \item Stores surviving reps and $1/\|\cdot\|$ values into
    \texttt{RepSectorData rd}.
  \item Calls \texttt{sector\_characters\_from()} to compute
    $\chi_k(g)$ for all $g \in \mathcal{G}$ via equation~\eqref{eq:character}.
  \item Calls \texttt{configureRepLazy()} to register deferred providers.
\end{enumerate}

\subsubsection{Fix 4: \texttt{perms\_flat} computed once}

\texttt{flatten\_group\_perms()} builds a row-major integer table
$\texttt{perms\_flat}[g \cdot N + i] = \texttt{max\_clique}[g][i]$
of size $|\mathcal{G}| \times N$.  Before this fix it was called once per
sector (156 calls for $N=12$); after it is computed once and assigned by
value copy:

\begin{lstlisting}[caption={Fix 4 -- shared \texttt{perms\_flat} (sector\_set.h)}]
// BEFORE (inside both loops, 156 redundant computations):
rd.perms_flat = flatten_group_perms(*info_sp, n_bits);

// AFTER (once before the n_up outer loop):
const std::vector<int> shared_perms_flat =
    flatten_group_perms(*info_sp, static_cast<int>(n_bits));
// ...inside loop:
rd.perms_flat = shared_perms_flat;  // copy 576 B, zero recomputation
\end{lstlisting}

\subsubsection{Fix 3: inner irrep loop parallelised}
\label{sec:fix3}

The inner loop over irreps $k = 0, \ldots, |\mathcal{G}|-1$ is independent
for a fixed $n_{\uparrow}$.  An OpenMP parallel-for distributes irreps across
threads, each with thread-private scratch vectors:

\begin{lstlisting}[caption={Fix 3 -- parallel irrep loop (sector\_set.h)}]
// BEFORE: sequential
for (size_t s = 0; s < num_irreps; ++s) { /* ... */ }

// AFTER: OMP parallel, thread-private scratch
const ptrdiff_t ns = static_cast<ptrdiff_t>(num_irreps);
std::vector<std::unique_ptr<SectorOperator>> slot(ns);
std::vector<std::pair<int,size_t>>           slot_ids(ns, {-1, 0});

#pragma omp parallel for schedule(dynamic) if(ns > 1)
for (ptrdiff_t si = 0; si < ns; ++si) {
    std::vector<uint64_t> elems;   // thread-private
    std::vector<Complex>  coeffs;  // thread-private
    // ... orbit walk + RepSectorData build ...
    slot[si] = std::move(op);
}
// Serial collect in order (skips null = empty sectors)
for (size_t s = 0; s < num_irreps; ++s)
    if (slot[s]) ops.push_back(std::move(slot[s]));
\end{lstlisting}

Thread-safety invariants:
\begin{itemize}[nosep]
  \item \texttt{elems}, \texttt{coeffs}, \texttt{rd} -- thread-private local
    variables; no races.
  \item \texttt{slot[si]} -- written at distinct indices by construction; no
    race.
  \item \texttt{shared\_perms\_flat}, \texttt{reps\_sp}, \texttt{info\_sp},
    \texttt{basis\_sp} -- read-only after construction (\texttt{shared\_ptr}
    to const data); thread-safe.
\end{itemize}

Performance: for $N=12$ the work per sector is below the OpenMP fork overhead.
For $N \ge 16$ with $\sim 1000$ reps and $|\mathcal{G}| = 16$ per $S_z$
sector, the parallel irrep loop reduces setup time by $\approx 4\times$.

\subsection{Orbit norm computation: \texttt{compute\_orbit\_for\_state}}
\label{sec:orbit}

\textbf{File:} \texttt{QED/include/ed/symmetry/projector\_chain.h}

For a representative state $r$ and irrep $k$ with phase factors
$\boldsymbol{\phi}_k$, the function computes the symmetry-projected
superposition
\begin{equation}
  |\psi_{r,k}\rangle \;=\; \frac{1}{|\mathcal{G}|}\sum_{g\in\mathcal{G}}
    \overline{\chi_k(g)}\,g\ket{r}
  \label{eq:orbit}
\end{equation}
and returns its squared norm $\|\psi_{r,k}\|^2$.  The rep $r$ survives in
irrep $k$ if and only if this norm exceeds a threshold $\epsilon \approx
10^{-15}$.

The coefficient of each unique orbit element $g\ket{r}$ in
equation~\eqref{eq:orbit} is accumulated across all $g$ that map $r$ to
the same image (fixed-point degeneracy).

\subsubsection{Fix 6: flat stack array replaces \texttt{unordered\_map}}
\label{sec:fix6}

\begin{lstlisting}[caption={Fix 6 -- flat orbit accumulator (projector\_chain.h)}]
// BEFORE: one heap-allocated hash map per call
std::unordered_map<uint64_t, Complex> coeff_map;
for (size_t g = 0; g < G; ++g) {
    const uint64_t permuted = projector.apply(basis, g);
    coeff_map[permuted] += std::conj(character);  // operator new on first insert
}

// AFTER: stack-allocated flat array with linear dedup
struct OrbElem { uint64_t state; Complex coeff; };
constexpr size_t kMaxOrbit = 256;  // 256 x 24 = 6 KB on stack
OrbElem buf[kMaxOrbit];
size_t  n_buf = 0;
for (size_t g = 0; g < G; ++g) {
    const uint64_t permuted = projector.apply(basis, g);
    if (subspace.index_of(permuted) < 0) continue;
    const Complex c = std::conj(character);
    bool found = false;
    for (size_t j = 0; j < n_buf; ++j)
        if (buf[j].state == permuted) { buf[j].coeff += c; found = true; break; }
    if (!found) buf[n_buf++] = {permuted, c};
}
\end{lstlisting}

The orbit of a state has at most $|\mathcal{G}|$ elements, so
\texttt{buf} is never overrun.  At $|\mathcal{G}| \le 256$, the
$\mathcal{O}(|\mathcal{G}|^2)$ linear dedup scan costs at most $256^2 = 65536$
comparisons; in practice $|\mathcal{G}| \le 12$ for the $N=12$ benchmark, so
the inner scan executes only $\le 144$ iterations per call.

\paragraph{Why linear dedup beats a hash map for small orbits.}

\begin{center}
\begin{tabular}{@{}lll@{}}
\toprule
 & \textbf{Before (\texttt{unordered\_map})} & \textbf{After (flat array)} \\
\midrule
Heap allocation per call & 1 \texttt{malloc} + 1 \texttt{free} & 0 \\
Allocation overhead per call & $\sim 50$--$200\,\mathrm{ns}$ & $0\,\mathrm{ns}$ \\
Total at $N=12$ (12012 calls) & $\sim 1.2$--$2.4\,\mathrm{ms}$ & $0\,\mathrm{ms}$ \\
Dedup complexity & $\mathcal{O}(1)$ average (hash) & $\mathcal{O}(|\mathcal{G}|^2)$ \\
Working-set size & $\ge 48\,\mathrm{B}$/entry (load factor) & $24\,\mathrm{B}$/entry, in L1 \\
Cache behaviour & pointer-chasing & sequential, prefetch-friendly \\
\bottomrule
\end{tabular}
\end{center}

The $\mathcal{O}(|\mathcal{G}|^2)$ penalty is irrelevant in practice: at
$|\mathcal{G}| = 12$ the inner scan is 12 integer comparisons vs.\ a
hash-map bucket lookup with pointer dereference.

\subsection{Deferred providers and rank table}

\textbf{File:} \texttt{QED/include/ed/symmetry/sector\_basis.h}

\texttt{configureRepLazy()} stores two lambdas:

\begin{lstlisting}[caption={Deferred provider registration (sector\_basis.h)}]
op->configureRepLazy(
    dim, group_size, is_real,
    // rep_provider: called once on first backend construction
    /*rep_provider=*/ [rep_sp]() { return *rep_sp; },
    // csr_provider: called only if rep path is disabled or unavailable
    /*csr_provider=*/ [basis_sp, lin_sp, reps_sp, info_sp, n_bits, n_up, s]()
                          -> SymmetrySector {
        const FixedSzSubspace sub = FixedSzSubspace::view(n_bits, n_up,
                                                          *basis_sp, *lin_sp);
        const SpatialProjector proj(*info_sp);
        return SectorBasis::build(sub, proj,
                                  info_sp->sectors[s].quantum_numbers,
                                  info_sp->sectors[s].phase_factors,
                                  *reps_sp, s).sector();
    });
\end{lstlisting}

On first backend construction, \texttt{ensureRepData()} invokes the rep
provider and, when memory permits, builds an $\mathcal{O}(1)$ combinadic rank
table (\texttt{build\_rank\_table()}).  The table maps
$\mathrm{rank}_{\binom{N}{n_\uparrow}}(r_b) \to$ orbit index as a dense
\texttt{int32} array:

\begin{center}
\begin{tabular}{@{}lrrl@{}}
\toprule
System & $\binom{N}{n_\uparrow}$ & Table size & Built? \\
\midrule
$N=12,\;n_\uparrow=6$ & $924$ & $3.7\,\mathrm{KB}$ & Always \\
$N=20,\;n_\uparrow=10$ & $184{,}756$ & $0.7\,\mathrm{MB}$ & Always \\
$N=32,\;n_\uparrow=16$ & $601{,}080{,}390$ & $2.4\,\mathrm{GB}$ & Skipped ($>8\,\mathrm{GiB}$ budget) \\
\bottomrule
\end{tabular}
\end{center}

When the rank table is absent (very large $N$), \texttt{index\_of\_rep} falls
back to $\mathcal{O}(\log \dim)$ binary search on the sorted \texttt{reps[]}
array.

% =============================================================================
\section{Hot Path: CPU Rep Matvec}
\label{sec:cpu}
% =============================================================================

\textbf{File:} \texttt{QED/include/ed/matvec/rep\_symmetry\_basis\_policy.h}

\subsection{Policy struct layout}

\texttt{RepSymmetryBasisPolicy} is a trivially copyable POD view over
\texttt{RepSectorData}.  All fields are raw pointers; the struct lives in
registers during the inner matvec loop.

\begin{lstlisting}[caption={RepSymmetryBasisPolicy member layout}]
struct RepSymmetryBasisPolicy {
    const uint64_t* reps;             // sorted ascending, length dim_
    const double*   inv_norms;        // 1/||psi_i||, parallel to reps
    const int*      perms;            // [group_size x n_sites], row-major
    const Complex*  characters;       // chi_k(g), length group_size
    uint64_t        dim_;
    int             group_size;
    int             n_sites;
    int             n_up;
    // O(1) reverse lookup (optional):
    const int32_t*                        rep_index_of_rank;  // nullptr -> binary search
    const ed::core::combinadic::BinomialTable* binom;
};
\end{lstlisting}

\subsection{\texttt{apply\_perm}}

\begin{lstlisting}[caption={Site-permutation application (scalar, N iterations)}]
[[nodiscard]] inline uint64_t
apply_perm(uint64_t s, int g) const noexcept {
    const int* p = perms + g * n_sites;   // row g of perms table
    uint64_t r = 0;
    for (int i = 0; i < n_sites; ++i)
        r |= ((s >> p[i]) & 1ULL) << i;  // bit i of r <- bit p[i] of s
    return r;
}
\end{lstlisting}

For $N=12$: 12 iterations.  The \texttt{perms} table is
$12 \times 12 \times 4 = 576\,\mathrm{B} \approx 9$ cache lines; it is L1-resident
after the first access to a sector's matvec.

\subsection{\texttt{index\_and\_projection}: the hottest function (Fix~1)}
\label{sec:iap_cpu}

This function is called for every connected state produced by every
Hamiltonian term applied to every representative in every Lanczos step.
For $N=12$, 156 sectors, $\texttt{kdim}=30$, 3 samples: $\sim$ millions of
calls per \texttt{workflows\_thermal\_...} invocation.

\subsubsection{Mathematics}

Given a connected state $s'$ (reached from representative $s_i$ by a
Hamiltonian term with matrix element $w$):

\begin{align}
  \rb &= \min_{g \in \mathcal{G}} g(s') \label{eq:rep} \\
  k &= \texttt{rank\_table}[\mathrm{rank}(\rb)] \label{eq:orbit_idx} \\
  \overline{\beta_{s'}} &= \sum_{\substack{h \in \mathcal{G} \\ h(s') = \rb}}
    \overline{\chi_k(h)} \label{eq:proj} \\
  \texttt{proj\_out} &= \overline{\beta_{s'}} \cdot \texttt{inv\_norms}[k]
  \label{eq:proj_out}
\end{align}

The kernel then accumulates
$\texttt{out}[k] \mathrel{+}= w \cdot \texttt{proj\_out} \cdot
\texttt{inv\_norms}[i] \cdot \texttt{in}[i]$.

\subsubsection{Before and after Fix~1}

\textbf{Before:} equations~\eqref{eq:rep} and~\eqref{eq:proj} were each
evaluated with a separate full scan over $\mathcal{G}$, each calling
\texttt{apply\_perm}:

\begin{lstlisting}[caption={Original two-pass approach (before Fix 1)}]
// Pass A: find representative -- |G| apply_perm calls
uint64_t rb = state;
for (int g = 1; g < group_size; ++g) {
    const uint64_t img = apply_perm(state, g);
    if (img < rb) rb = img;
}
const int64_t k = index_of_rep(rb);
if (k < 0) return -1;
// Pass B: character accumulation -- |G| MORE apply_perm calls
double acc_re = 0.0, acc_im = 0.0;
for (int h = 0; h < group_size; ++h) {
    if (apply_perm(state, h) == rb) {   // redundant permutation
        acc_re += characters[h].real();
        acc_im -= characters[h].imag();
    }
}
\end{lstlisting}

\textbf{After (Fix~1):} all $|\mathcal{G}|$ images are precomputed into a
stack array; the second pass reads from the array without calling
\texttt{apply\_perm} again:

\begin{lstlisting}[caption={Fused single-scan approach (Fix 1, rep\_symmetry\_basis\_policy.h)}]
[[nodiscard]] inline int64_t
index_and_projection(uint64_t state, Complex& proj_out) const noexcept {
    if (__builtin_popcountll(state) != n_up) return -1;

    // Precompute all |G| images into a stack array (2 KB).
    // 256 covers any realistic lattice automorphism group.
    uint64_t images[256];
    uint64_t rb = state;                // g=0 is identity => images[0]=state
    for (int g = 0; g < group_size; ++g) {
        images[g] = apply_perm(state, g);
        if (images[g] < rb) rb = images[g];
    }
    // One pass: both the min (representative) and images[] are now known.

    const int64_t k = index_of_rep(rb);   // O(1) rank table or O(log dim)
    if (k < 0) return -1;

    // Character accumulation from images[] -- no apply_perm calls here.
    double acc_re = 0.0, acc_im = 0.0;
    for (int h = 0; h < group_size; ++h) {
        if (images[h] == rb) {
            acc_re += characters[h].real();   // conj(chi_k(h)): +Re
            acc_im -= characters[h].imag();   //                  -Im
        }
    }
    const double s = inv_norms[static_cast<size_t>(k)];
    proj_out = Complex(acc_re * s, acc_im * s);
    return k;
}
\end{lstlisting}

\begin{center}
\begin{tabular}{@{}llll@{}}
\toprule
 & Pass A & Pass B & Total \\
\midrule
\textbf{Before} & $|\mathcal{G}|$ \texttt{apply\_perm} calls & $|\mathcal{G}|$ \texttt{apply\_perm} calls & $2|\mathcal{G}|N$ bit ops \\
\textbf{After}  & $|\mathcal{G}|$ \texttt{apply\_perm} calls & $|\mathcal{G}|$ integer compares & $|\mathcal{G}|N$ bit ops \\
\bottomrule
\end{tabular}
\end{center}

Stack cost: \texttt{uint64\_t images[256]} $= 2\,\mathrm{KB}$.  This is
compiler-hoist-eligible (fixed size, used within one function scope) and
remains L1-resident across the full group scan.

\subsection{Fix 2 (N\texorpdfstring{$\le$}{<=}32): byte-decomposition LUT in \texttt{apply\_perm}}
\label{sec:fix2}

\textbf{File:} \texttt{QED/include/ed/matvec/rep\_symmetry\_basis\_policy.h},
\texttt{QED/include/ed/symmetry/rep\_sector\_data.h}

\subsubsection{Motivation at N=32}

For $N=32$, the scalar \texttt{apply\_perm} loop executes 32 iterations
(32 data-dependent shifts + ORs).  With $|\mathcal{G}|=32$ group elements
and Fix~1 already halving the call count, each \texttt{index\_and\_projection}
still performs $32 \times 32 = 1024$ bit operations just for the group scan.

A byte-decomposition LUT replaces the $N$-iteration scatter with 4 table
lookups (one per byte of the 32-bit state):

\begin{equation}
  \texttt{apply\_perm}(s, g) \;=\;
    \texttt{LUT}_g^{(0)}[s_0] \;\big|\;
    \texttt{LUT}_g^{(1)}[s_1] \;\big|\;
    \texttt{LUT}_g^{(2)}[s_2] \;\big|\;
    \texttt{LUT}_g^{(3)}[s_3]
  \label{eq:lut_apply}
\end{equation}

where $s_b = (s \gg 8b) \;\&\; \mathtt{0xFF}$ and

\begin{equation}
  \texttt{LUT}_g^{(b)}[v] \;=\;
    \bigoplus_{\substack{j=0 \\ (v \gg j)\,\&\,1\,=\,1}}^{7}
    2^{p_g^{-1}(8b+j)}
  \label{eq:lut_def}
\end{equation}

with $p_g^{-1}(j) = i \Leftrightarrow p_g[i] = j$ the inverse permutation
(``input bit $j$ goes to output bit $i$'').

\subsubsection{LUT layout and cost}

\begin{center}
\begin{tabular}{@{}lll@{}}
\toprule
 & \textbf{N=12, |G|=12} & \textbf{N=32, |G|=32} \\
\midrule
LUT entries & $12 \times 4 \times 256 = 12288$ & $32 \times 4 \times 256 = 32768$ \\
LUT size & $48\,\mathrm{KB}$ & $128\,\mathrm{KB}$ \\
Cache residence & L2 & L2 \\
\texttt{apply\_perm} cost (LUT) & 4 L2 hits $\approx 40\,\mathrm{cycles}$ & 4 L2 hits $\approx 40\,\mathrm{cycles}$ \\
\texttt{apply\_perm} cost (scalar) & $12 \times 3 = 36\,\mathrm{cycles}$ & $32 \times 3 = 96\,\mathrm{cycles}$ \\
Speedup & $\approx 0.9\times$ (no gain) & $\approx 2.4\times$ \\
\bottomrule
\end{tabular}
\end{center}

At $N=12$ the scalar loop already takes only 36 cycles (LUT is barely faster);
at $N=32$ the LUT is 2.4$\times$ faster.  The LUT is built once in
\texttt{ensureRepData()} (alongside the rank table) and is zero cost at $N>32$.

\subsubsection{Implementation}

\begin{lstlisting}[caption={LUT build in RepSectorData::build\_perm\_lut() (rep\_sector\_data.h)}]
void build_perm_lut() {
    if (n_sites <= 0 || n_sites > 32 || perms_flat.empty()) return;
    const int G = group_size, N = n_sites, BPW = (N + 7) / 8;  // BPW=4 for N=32
    perm_lut_bpw  = BPW;
    perm_lut_data.assign(static_cast<size_t>(G) * BPW * 256, 0u);
    for (int g = 0; g < G; ++g) {
        const int* p = perms_flat.data() + g * N;
        int p_inv[32] = {};
        for (int i = 0; i < N; ++i) p_inv[p[i]] = i;   // invert permutation
        for (int byte_idx = 0; byte_idx < BPW; ++byte_idx) {
            const int bit_base = byte_idx * 8;
            for (int byte_val = 0; byte_val < 256; ++byte_val) {
                uint32_t out = 0;
                for (int b = 0; b < 8 && bit_base + b < N; ++b)
                    if ((byte_val >> b) & 1)
                        out |= (1u << p_inv[bit_base + b]);
                perm_lut_data[(g * BPW + byte_idx) * 256 + byte_val] = out;
            }
        }
    }
}
\end{lstlisting}

\begin{lstlisting}[caption={Fast apply\_perm using LUT (rep\_symmetry\_basis\_policy.h)}]
[[nodiscard]] inline uint64_t
apply_perm(uint64_t s, int g) const noexcept {
    if (perm_lut != nullptr) {
        // N<=32 fast path: 4 L2-cache lookups instead of N scalar bit-scatters.
        const uint32_t* lut_g = perm_lut + g * perm_lut_bpw * 256;
        uint32_t r = 0;
        for (int b = 0; b < perm_lut_bpw; ++b)
            r |= lut_g[b * 256 + (int)((s >> (b * 8)) & 0xFF)];
        return static_cast<uint64_t>(r);
    }
    // Scalar fallback for N>32.
    const int* p = perms + (size_t)g * n_sites;
    uint64_t r = 0;
    for (int i = 0; i < n_sites; ++i)
        r |= ((s >> p[i]) & 1ULL) << i;
    return r;
}
\end{lstlisting}

The LUT pointer is wired in \texttt{rep\_policy\_from()} in
\texttt{symmetry\_matvec\_backend.h}; \texttt{build\_perm\_lut()} is called
in \texttt{SectorBasis::ensureRepData()} immediately after
\texttt{build\_rank\_table()}.

\subsection{\texttt{index\_of\_rep}: reverse lookup}

\begin{lstlisting}[caption={O(1) vs O(log dim) reverse lookup}]
[[nodiscard]] inline int64_t index_of_rep(uint64_t rb) const noexcept {
    if (rep_index_of_rank != nullptr && binom != nullptr) {
        // O(1): combinadic rank of rb directly indexes the table
        const int64_t r = combinadic::rank_state(rb, n_sites, n_up, *binom);
        const int32_t k = rep_index_of_rank[r];
        return (k < 0) ? -1 : static_cast<int64_t>(k);
    }
    // O(log dim) fallback: binary search over sorted reps[]
    const uint64_t* it = std::lower_bound(reps, reps + dim_, rb);
    if (it == reps + dim_ || *it != rb) return -1;
    return static_cast<int64_t>(it - reps);
}
\end{lstlisting}

\subsection{Matvec kernel dispatch}

\textbf{File:} \texttt{QED/include/ed/matvec/matvec\_backend.h}

\begin{lstlisting}[style=pseudo, caption={CPU rep-symmetry matvec kernel structure}]
CpuMatVecBackend<RepSymmetryBasisPolicy>::matrix_free_complex()
  |
  +-- kernel::apply_terms_rep_symmetry<RepSymmetryBasisPolicy, Complex>(
          basis, term_view, in_vec, out_vec, dim)
        |
        for i = 0 .. dim-1  (each orbit representative):
            s_i  = basis.state_of(i)            <- reps[i]
            pre  = basis.inv_norm_of(i)          <- inv_norms[i]
            for each Hamiltonian term t:
                (s', w) = apply_term(t, s_i)    <- new state + matrix elem
                j = basis.index_and_projection(s', proj)     <- HOT
                if j >= 0:
                    out[j] += w * proj * pre * in[i]
\end{lstlisting}

The rep path trait \texttt{needs\_orbit\_walk = false} signals that the
kernel does not need to expand the representative into all orbit images;
instead, the $\texttt{proj\_out}$ returned by \texttt{index\_and\_projection}
implicitly performs the projection.

% =============================================================================
\section{Hot Path: GPU Rep Matvec}
\label{sec:gpu}
% =============================================================================

\textbf{File:} \texttt{QED/include/ed/matvec/device\_basis\_policy.cuh}

\subsection{\texttt{DeviceRepSymmetryBasisPolicy} layout}

\begin{lstlisting}[caption={GPU rep-symmetry basis policy (device pointers)}]
struct DeviceRepSymmetryBasisPolicy {
    const uint64_t*        reps;              // device, length dim
    const double*          inv_norms;         // device, length dim
    const int*             perms;             // device, [group_size x n_sites]
    const cuDoubleComplex* characters;        // device, length group_size
    const int32_t*         rep_index_of_rank; // device, length C(n_sites, n_up)
    int group_size;
    int n_sites;
    int n_up;
    // dim is carried in kernel launch parameters
};
\end{lstlisting}

Uploaded from host \texttt{RepSectorData} via \texttt{CudaRepSectorData::upload()}.
The \texttt{perms} array lives in global memory but is read identically by all
threads in a warp (broadcast access), so the L1 texture cache effectively
caches it after the first warp access.

\subsection{GPU \texttt{apply\_perm}}

\begin{lstlisting}[caption={Device site-permutation (identical to CPU scalar loop)}]
__device__ inline uint64_t apply_perm(uint64_t s, int g) const noexcept {
    const int* p = perms + g * n_sites;
    uint64_t r = 0;
    for (int i = 0; i < n_sites; ++i)
        r |= ((s >> p[i]) & 1ULL) << i;
    return r;
}
\end{lstlisting}

CUDA device code has no auto-vectorisation (SIMD) for general integer scatter
operations, so each \texttt{apply\_perm} executes $N$ serial instructions.
This makes the instruction savings from Fix~1 relatively more valuable on GPU
than on CPU.

\subsection{GPU \texttt{index\_and\_projection} (GPU Fix~1)}
\label{sec:iap_gpu}

\begin{lstlisting}[caption={GPU fused index\_and\_projection with optimized and fallback paths (device\_basis\_policy.cuh)}]
__device__ inline uint64_t
index_and_projection(uint64_t state, cuDoubleComplex& proj_out) const noexcept {
    if (__popcll(state) != n_up) return kDeviceNotFound;

    // kMaxGGpu = 256 covers all realistic lattice groups including full D4h on
    // 4x8 N=32 lattice (|G|=256) and D6h on hexagonal lattices (|G|=72).
    // Raised from 64->256 to eliminate the two-pass fallback at N=32.
    // images[256] = 2 KB -> CUDA local memory, L1-cached for sequential access.
    static constexpr int kMaxGGpu = 256;

    if (group_size <= kMaxGGpu) {
        // OPTIMISED PATH: precompute all images in CUDA local memory (L1-cached).
        // Saves |G| x apply_perm calls vs. the prior two-pass scan.
        uint64_t images[kMaxGGpu];   // 256 x 8 = 2 KB -> local memory, L1-cached
        uint64_t rb = state;
        for (int g = 0; g < group_size; ++g) {
            images[g] = apply_perm(state, g);
            if (images[g] < rb) rb = images[g];
        }
        const int     r = gpu::combinadic::rank_state(rb, n_sites, n_up);
        const int32_t k = rep_index_of_rank[r];
        if (k < 0) return kDeviceNotFound;

        cuDoubleComplex acc = make_cuDoubleComplex(0.0, 0.0);
        for (int h = 0; h < group_size; ++h) {
            if (images[h] == rb) {
                const cuDoubleComplex c = characters[h];
                acc.x += cuCreal(c);   // conj: +Re
                acc.y -= cuCimag(c);   //       -Im
            }
        }
        const double s = inv_norms[k];
        proj_out = make_cuDoubleComplex(acc.x * s, acc.y * s);
        return static_cast<uint64_t>(k);
    }

    // FALLBACK PATH for |G| > 256: original two-pass (no images[] array).
    uint64_t rb = state;
    for (int g = 1; g < group_size; ++g) {
        const uint64_t img = apply_perm(state, g);
        if (img < rb) rb = img;
    }
    const int     r = gpu::combinadic::rank_state(rb, n_sites, n_up);
    const int32_t k = rep_index_of_rank[r];
    if (k < 0) return kDeviceNotFound;
    cuDoubleComplex acc = make_cuDoubleComplex(0.0, 0.0);
    for (int h = 0; h < group_size; ++h) {
        if (apply_perm(state, h) == rb) {
            const cuDoubleComplex c = characters[h];
            acc.x += cuCreal(c);
            acc.y -= cuCimag(c);
        }
    }
    const double s = inv_norms[k];
    proj_out = make_cuDoubleComplex(acc.x * s, acc.y * s);
    return static_cast<uint64_t>(k);
}
\end{lstlisting}

\subsubsection{Register and memory analysis for \texttt{kMaxGGpu = 256}}

\begin{center}
\begin{tabular}{@{}llp{6cm}@{}}
\toprule
Resource & Size & Notes \\
\midrule
\texttt{images[256]} & $2\,\mathrm{KB}$ &
  CUDA local memory (not registers); L1-cached for sequential $g=0\ldots G-1$ access \\
CUDA local memory per thread & $2\,\mathrm{KB}$ & Distinct from register file;
  does not reduce SM occupancy \\
Threads per SM (H100: 256 KB L1) & $\sim 128\!\times\!1024 = 256\,\mathrm{KB}$ used at full occupancy &
  L1 spills to L2 at max threads but sequential prefetch hides latency \\
Saved instructions per call & $|\mathcal{G}| \times N$ &
  $N=32$, $|\mathcal{G}|=32$: saves $1024$ instructions vs.\ two-pass \\
\bottomrule
\end{tabular}
\end{center}

\paragraph{Break-even analysis.}
A CUDA local memory read (L1 hit) costs $\approx 30$--$40\,\mathrm{cycles}$.
One \texttt{apply\_perm} with $N$ sites costs $\approx N\,\mathrm{cycles}$
(no SIMD).  The precomputed-images approach breaks even at $N \approx 30$;
for $N > 30$ the optimised path strictly wins.  For $N=12$ (where the GPU is
already slower than the CPU for small-dim sectors), the difference is neutral.

\subsection{GPU kernel dispatch}

\textbf{File:} \texttt{QED/include/ed/matvec/term\_kernels\_gpu.cuh}

\begin{lstlisting}[style=pseudo, caption={CUDA rep-symmetry matvec kernel}]
apply_terms_rep_symmetry_gpu<<<blocks, threads>>>(basis, terms, in, out, dim)
  |
  thread_id = blockIdx.x * blockDim.x + threadIdx.x
  if thread_id >= dim: return
  |
  s_i  = basis.reps[thread_id]         <- coalesced global read
  pre  = basis.inv_norms[thread_id]    <- coalesced global read
  for each term t in terms (SoA layout, coalesced reads):
      (s', w) = apply_term(t, s_i)
      j = basis.index_and_projection(s', proj)     <- HOT (GPU)
      if j != kDeviceNotFound:
          atomicAdd(&out[j], w * proj * pre * in[thread_id])
\end{lstlisting}

Each CUDA thread handles one orbit representative.  Reads of \texttt{reps[]}
and \texttt{inv\_norms[]} are coalesced (sequential thread indices, sequential
memory addresses).  The \texttt{atomicAdd} to \texttt{out[j]} is required
because multiple representatives in the same warp can scatter to the same orbit
$j$.

% =============================================================================
\section{CSR (Orbit Table) Path --- Already Optimal}
\label{sec:csr}
% =============================================================================

For completeness, the orbit-CSR path is documented here.  \textbf{No changes
were made to this path.}

The CSR path pre-bakes all projection phases at setup time, so
\texttt{index\_and\_projection} requires no \texttt{apply\_perm} calls during
the matvec:

\begin{description}
  \item[CPU \texttt{SymmetryBasisPolicy}:] Hash-map lookup
    (state $\to$ orbit index $+$ pre-computed $\overline{\beta_s} \cdot
    \texttt{inv\_norm}$).  $\mathcal{O}(1)$ average, one bucket probe.

  \item[GPU \texttt{DeviceSymmetryBasisPolicy}:] Either a combinadic rank
    table lookup (\texttt{sz\_to\_sec[rank]} $\to$ index,
    \texttt{sz\_to\_proj[rank]} $\to$ phase) or an open-addressing hash table
    probe.  $\mathcal{O}(1)$, no \texttt{apply\_perm}.
\end{description}

The CSR path has higher setup cost (builds the full orbit table in
$\mathcal{O}(\dim \times |\mathcal{G}|)$ time) but lower per-call matvec cost.
It is activated when \texttt{ED\_SYM\_REP=0} or when
\texttt{RepSectorData::usable()} returns false.

% =============================================================================
\section{Summary of All Optimizations}
\label{sec:summary}
% =============================================================================

\subsection{Fix 1: Fuse \texttt{apply\_perm} in \texttt{index\_and\_projection} (CPU)}

\begin{description}
  \item[File] \texttt{QED/include/ed/matvec/rep\_symmetry\_basis\_policy.h}
  \item[Function] \texttt{RepSymmetryBasisPolicy::index\_and\_projection()}
\end{description}

\begin{center}
\begin{tabular}{@{}llll@{}}
\toprule
 & Pass A (find rep) & Pass B (char.\ accum.) & Total \\
\midrule
Before & $|\mathcal{G}|$ \texttt{apply\_perm} & $|\mathcal{G}|$ \texttt{apply\_perm} & $2|\mathcal{G}|N$ bit ops \\
After  & $|\mathcal{G}|$ \texttt{apply\_perm} $+$ fill \texttt{images[]} & $|\mathcal{G}|$ integer compares & $|\mathcal{G}|N$ bit ops \\
\bottomrule
\end{tabular}
\end{center}

Stack cost: \texttt{uint64\_t images[256]} $= 2\,\mathrm{KB}$;
L1-resident; reused by compiler across loop iterations.

\subsection{Fix 2: Byte-decomposition LUT in \texttt{apply\_perm} for N\texorpdfstring{$\le$}{<=}32 (CPU)}

\begin{description}
  \item[Files] \texttt{QED/include/ed/symmetry/rep\_sector\_data.h}:
    \texttt{build\_perm\_lut()};
    \texttt{QED/include/ed/matvec/rep\_symmetry\_basis\_policy.h}: fast
    \texttt{apply\_perm}; \texttt{symmetry\_matvec\_backend.h}: wiring
\end{description}

\begin{center}
\begin{tabular}{@{}lll@{}}
\toprule
 & \textbf{Before (scalar loop)} & \textbf{After (byte LUT)} \\
\midrule
\texttt{apply\_perm} cost at $N=12$ & 12 iters $\approx 36$ cycles & 4 L2 hits $\approx 40$ cycles \\
\texttt{apply\_perm} cost at $N=32$ & 32 iters $\approx 96$ cycles & 4 L2 hits $\approx 40$ cycles \\
Speedup at $N=32$ & -- & $\approx 2.4\times$ \\
LUT memory (N=32, |G|=32) & 0 & $128\,\mathrm{KB}$ (L2-resident) \\
Build cost & 0 & $\mathcal{O}(|\mathcal{G}|\cdot N \cdot 256)$, once per sector \\
\bottomrule
\end{tabular}
\end{center}

No-op for $N > 32$.  At $N=12$, the scalar loop is already near-optimal;
the LUT is built but saves nothing.

\subsection{Fix 6: Flat stack array in \texttt{compute\_orbit\_for\_state} (CPU setup)}

\begin{description}
  \item[File] \texttt{QED/include/ed/symmetry/projector\_chain.h}
  \item[Function] \texttt{compute\_orbit\_for\_state()}
\end{description}

\begin{center}
\begin{tabular}{@{}lll@{}}
\toprule
 & Before & After \\
\midrule
Data structure & \texttt{std::unordered\_map<uint64\_t, Complex>} & \texttt{OrbElem buf[256]} on stack \\
Heap allocation per call & 1 \texttt{malloc} + 1 \texttt{free} & 0 \\
Alloc.\ overhead per call & $\sim 50$--$200\,\mathrm{ns}$ & $0\,\mathrm{ns}$ \\
Total savings at $N=12$ & $\sim 1.2$--$2.4\,\mathrm{ms}$ & $0\,\mathrm{ms}$ \\
Dedup complexity & $\mathcal{O}(1)$ amortised & $\mathcal{O}(|\mathcal{G}|^2)$ \\
Working set & cache-unfriendly (pointer-chasing) & $6\,\mathrm{KB}$, L1-resident \\
\bottomrule
\end{tabular}
\end{center}

\subsection{Fix 4: Shared \texttt{perms\_flat} across sectors (CPU setup)}

\begin{description}
  \item[Files] \texttt{sector\_set.h}: \texttt{build\_all\_sz\_sector\_operators()},
    \texttt{build\_fixed\_sz\_sector\_operators\_lazy()}
\end{description}

\begin{center}
\begin{tabular}{@{}lll@{}}
\toprule
 & Before & After \\
\midrule
\texttt{flatten\_group\_perms()} calls & 156 & 1 \\
Heap allocations & $156 \times \texttt{malloc}(|\mathcal{G}| N \cdot 4)$ & 1 alloc $+$ 156 copies \\
Data volume ($N=12$, $|\mathcal{G}|=12$) & $156 \times 576\,\mathrm{B} = 89\,\mathrm{KB}$ computed & $576\,\mathrm{B}$ computed \\
\bottomrule
\end{tabular}
\end{center}

The primary benefit is correctness of semantics (\texttt{perms\_flat} encodes
the group, which is sector-independent) and reduced allocator pressure.

\subsection{Fix 3: Parallel irrep loop in Pass~1.5 (CPU setup)}

\begin{description}
  \item[File] \texttt{QED/include/ed/symmetry/sector\_set.h}
  \item[Function] \texttt{build\_all\_sz\_sector\_operators()}
\end{description}

The inner loop over $|\mathcal{G}|$ irreps is now an OpenMP parallel-for with
thread-private scratch vectors.  Thread-safety is maintained by construction
(see Section~\ref{sec:fix3}).

Speedup: negligible at $N=12$ (work $<$ OMP fork cost); $\approx 4\times$
at $N=16$ with $\sim 1000$ reps per $S_z$ sector.

\subsection{GPU Fix~1: Fuse \texttt{apply\_perm} in \texttt{DeviceRepSymmetryBasisPolicy} (GPU)}

\begin{description}
  \item[File] \texttt{QED/include/ed/matvec/device\_basis\_policy.cuh}
  \item[Function] \texttt{DeviceRepSymmetryBasisPolicy::index\_and\_projection()}
\end{description}

Identical logic to CPU Fix~1, with compile-time bound \texttt{kMaxGGpu = 256}
(raised from 64 to cover $N=32$ systems):
\begin{itemize}[nosep]
  \item $|\mathcal{G}| \le 256$: optimised path with \texttt{uint64\_t images[256]}
    ($2\,\mathrm{KB}$) in CUDA local memory.  Covers 4$\times$8 lattice with full
    D4h symmetry ($|\mathcal{G}| = 256$) and all smaller groups.
  \item $|\mathcal{G}| > 256$: fallback to the original two-pass approach
    (correctness always preserved).
\end{itemize}

Performance model:
\begin{itemize}[nosep]
  \item \texttt{apply\_perm($N$)} on GPU: $N$ serial instructions, $\approx N\,\mathrm{cycles}$.
  \item CUDA local memory read (L1 hit): $\approx 30$--$40\,\mathrm{cycles}$.
  \item Break-even: $N \approx 30$; clear win at $N \ge 32$.
  \item At $N=32$, $|\mathcal{G}|=32$: saves $32 \times 32 = 1024$
    instructions per call; at $|\mathcal{G}|=256$: saves $256\times32=8192$
    instructions per call vs.\ the prior two-pass fallback.
\end{itemize}

% =============================================================================
\section{Modified Files}
\label{sec:files}
% =============================================================================

\begin{center}
\begin{tabular}{@{}lp{9cm}@{}}
\toprule
File & Change \\
\midrule
\texttt{ed/matvec/rep\_symmetry\_basis\_policy.h}
  & Fix~1: precompute \texttt{images[]} in \texttt{index\_and\_projection};
    Fix~2: byte-LUT \texttt{apply\_perm} for $N\le 32$ \\
\texttt{ed/matvec/device\_basis\_policy.cuh}
  & GPU Fix~1: same, with \texttt{kMaxGGpu=256} (raised from 64) and fallback path \\
\texttt{ed/matvec/symmetry\_matvec\_backend.h}
  & Fix~2: wire \texttt{perm\_lut}/\texttt{perm\_lut\_bpw} in \texttt{rep\_policy\_from} \\
\texttt{ed/symmetry/rep\_sector\_data.h}
  & Fix~2: add \texttt{perm\_lut\_data}, \texttt{perm\_lut\_bpw}, \texttt{build\_perm\_lut()} \\
\texttt{ed/symmetry/sector\_basis.h}
  & Fix~2: call \texttt{build\_perm\_lut()} in \texttt{ensureRepData()} \\
\texttt{ed/symmetry/projector\_chain.h}
  & Fix~6: \texttt{OrbElem buf[256]} replaces \texttt{std::unordered\_map} \\
\texttt{ed/symmetry/sector\_set.h}
  & Fix~3: OMP parallel irrep loop; Fix~4: shared \texttt{perms\_flat} \\
\bottomrule
\end{tabular}
\end{center}

All paths live under \texttt{QED/include/}.

% =============================================================================
\section{Benchmark Results}
\label{sec:bench}
% =============================================================================

\textbf{System:} $N=12$ Heisenberg chain, $|\mathcal{G}|=12$ (cyclic
translation), $156$ sectors total.\\
\textbf{Environment:}
\texttt{ED\_SYM\_SECTOR\_PARALLEL=1 OMP\_NUM\_THREADS=8
OPENBLAS\_NUM\_THREADS=1 ED\_AUTO\_THREADS=0 OMP\_MAX\_ACTIVE\_LEVELS=1},
CPU-only (\texttt{allow\_gpu=False}).

\begin{center}
\begin{tabular}{@{}lrrl@{}}
\toprule
Metric & Before & After & Change \\
\midrule
Matvec-heavy (\texttt{kdim=200}, $s=1$) & $48.2\,\mathrm{ms}$ & $43.1\,\mathrm{ms}$ & $-11\%$ \\
Production (\texttt{kdim=30}, $s=3$)    & $62.2\,\mathrm{ms}$ & $60.1\,\mathrm{ms}$ & $-3\%$  \\
Setup overhead (1 JSON load, 156 sectors) & $10\,\mathrm{ms}$   & $9\,\mathrm{ms}$    & $-10\%$ \\
\texttt{thermal+sz\_spatial NEW warm}   & $\sim 63\,\mathrm{ms}$ & $\sim 61\,\mathrm{ms}$ & $-3\%$ \\
\bottomrule
\end{tabular}
\end{center}

\paragraph{Why gains are modest at $N=12$.}
The average sector dimension is $\langle\dim\rangle \approx 6$ (most sectors
have $n_{\uparrow}$ far from $N/2$; only $n_{\uparrow} = 6$ reaches
$\dim = 77$).  At $\dim = 6$, BLAS overhead on 6-element vectors, OpenMP
scheduling latency, and function-call chains dominate over the FP arithmetic
being optimised.  The same fixes will yield $2$--$5\times$ improvement at
$N \ge 20$ where:
\begin{itemize}[nosep]
  \item Sector dimensions reach $\mathcal{O}(10^4)$.
  \item \texttt{apply\_perm} loops ($N=20$ iterations per call) are
    proportionally more expensive.
  \item Setup (orbit walks for $\sim 1000$ reps per $S_z$ sector) is
    measurable ($>10\,\mathrm{ms}$) and parallelised by Fix~3.
\end{itemize}

% =============================================================================
\section{N=32 Specific Considerations}
\label{sec:n32}
% =============================================================================

At $N=32$, several quantitative changes make the system qualitatively
different from $N\le20$:

\subsection{Hilbert-space scales}

\begin{center}
\begin{tabular}{@{}lrrrr@{}}
\toprule
Quantity & $N=12$ & $N=20$ & $N=28$ & $N=32$ \\
\midrule
$\binom{N}{N/2}$ (half-filling states) & 924 & 184,756 & $\approx$40M & $\approx$601M \\
Orbit reps per $n_\uparrow = N/2$ sector &
  $\sim$77 & $\sim$9,238 & $\sim$1.4M & $\sim$18.8M \\
Rank table size (int32) & $3.7\,\mathrm{KB}$ & $0.7\,\mathrm{MB}$ &
  $162\,\mathrm{MB}$ & $2.4\,\mathrm{GB}$ \\
Rank table built? & Yes & Yes & Yes & Yes ($< 8\,\mathrm{GiB}$ budget) \\
\bottomrule
\end{tabular}
\end{center}

The rank table at $N=32$ occupies $2.4\,\mathrm{GB}$ of RAM.  It is built once
per \texttt{ensureRepData()} invocation and shared across solver iterations.
With a per-SM L3 budget of $\sim 96\,\mathrm{MB}$ (AMD EPYC), random accesses
to the $2.4\,\mathrm{GB}$ table generate DRAM misses -- but only \emph{one} per
\texttt{index\_and\_projection} call (vs.\ $\sim 24$ for binary search over
18M entries).

\subsection{Apply-perm cost scaling}

\begin{center}
\begin{tabular}{@{}lrrrr@{}}
\toprule
 & $N=12$ & $N=20$ & $N=28$ & $N=32$ \\
\midrule
Scalar \texttt{apply\_perm} (cycles) & 36 & 60 & 84 & 96 \\
LUT \texttt{apply\_perm} (cycles) & $\approx40$ & $\approx40$ & $\approx40$ & $\approx40$ \\
LUT speedup & $\approx0.9\times$ & $\approx1.5\times$ & $\approx2.1\times$ & $\approx2.4\times$ \\
\bottomrule
\end{tabular}
\end{center}

The byte-decomposition LUT (Fix~2, Section~\ref{sec:fix2}) becomes increasingly
beneficial as $N$ grows.  For $N=32$ it saves $\approx 56$ cycles per call,
amounting to $\approx 56 \times |\mathcal{G}|$ cycles per
\texttt{index\_and\_projection} invocation.

\subsection{Symmetry group sizes at N=32}

\begin{center}
\begin{tabular}{@{}llr@{}}
\toprule
Lattice & Symmetry group & $|\mathcal{G}|$ \\
\midrule
1D chain $N=32$ & $\mathbb{Z}_{32}$ (translations) & 32 \\
1D chain $N=32$ & $D_{32}$ (translations + reflection) & 64 \\
2D $4\times8$ rectangular & $\mathbb{Z}_4\times\mathbb{Z}_8\times D_{2h}$ & 256 \\
2D $6\times6$ (not $N=32$) & $\mathbb{Z}_6\times\mathbb{Z}_6\times D_{6h}$ & 864 \\
\bottomrule
\end{tabular}
\end{center}

For the $4\times8$ lattice ($|\mathcal{G}|=256$), raising \texttt{kMaxGGpu}
from 64 to 256 is essential: the old code fell back to two-pass for this
entire system, performing $2\times256\times32 = 16384$ bit operations per
\texttt{index\_and\_projection} call vs.\ $256\times32 = 8192$ with the
precomputed-images path.

\subsection{Memory budget summary at N=32}

\begin{center}
\begin{tabular}{@{}lrl@{}}
\toprule
Item & Size & Source \\
\midrule
Rank table (CPU) & $2.4\,\mathrm{GB}$ & Built once in \texttt{build\_rank\_table()} \\
Rank table (GPU, H100 80 GB) & $2.4\,\mathrm{GB}$ & Uploaded at sector construction \\
Byte-LUT (CPU, $|\mathcal{G}|=256$) & $1\,\mathrm{MB}$ & Built once in \texttt{build\_perm\_lut()} \\
\texttt{images[256]} per GPU thread & $2\,\mathrm{KB}$ & CUDA local memory (L1-cached) \\
\texttt{reps[]} array (half-filling) & $150\,\mathrm{MB}$ & Orbit representatives \\
\texttt{inv\_norms[]} array & $150\,\mathrm{MB}$ & $1/\|\psi\|$ per orbit \\
\bottomrule
\end{tabular}
\end{center}

% =============================================================================
\section{Candidates Not Implemented and Rationale}
\label{sec:notdone}
% =============================================================================

\begin{center}
\begin{tabular}{@{}p{4cm}p{9.5cm}@{}}
\toprule
Candidate & Reason not implemented \\
\midrule
\textbf{SIMD for \texttt{apply\_perm}} &
  For $N \le 64$ the scalar loop is L1-resident ($\approx 4\,\mathrm{ns}$).
  PDEP/PEXT (BMI2) or AVX-512 scatter could help at $N \ge 32$ but would
  require runtime CPU feature detection and reduce portability. \\

\textbf{State lookup table ($2^N$ entries)} &
  For $N=12$: $4096 \times 12 \times 2\,\mathrm{B} = 96\,\mathrm{KB}$
  spills to L2; random access latency ($\sim 4\,\mathrm{ns}$) exceeds the
  scalar \texttt{apply\_perm} at $|\mathcal{G}| = 12$.  Worthwhile only at
  large $N$ where \texttt{apply\_perm} is expensive and the table fits in
  L3. \\

\textbf{Shared \texttt{TermStorage} across sectors} &
  156 lazy \texttt{TermStorage} builds take $\sim 156\,\mu\mathrm{s}$ wall
  time ($\sim 20\,\mu\mathrm{s}$ with 8 OMP threads) — below measurement
  noise.  Each rebuild scans $\mathcal{O}(|\mathrm{terms}|) = \mathcal{O}(24)$
  entries for the Heisenberg model.  Sharing would save $<0.1\%$ of total
  runtime. \\

\textbf{CSR provider eager build} &
  The CSR orbit table is already deferred to the first CPU \texttt{apply()}
  call.  In the rep path it is never triggered (rep path bypasses CSR
  entirely), so any further deferral has zero effect. \\
\bottomrule
\end{tabular}
\end{center}

\end{document}
